% Your comments for the rook-ending positions. Format: notes/README.md.

% Diagram 39 of the 1500-1599 book: Majstrova - Petruzalek
@ 8/8/4k3/8/8/1pr5/1R6/4K3 w - - 2 67 | Rd2?
Do not improvise the setup!\index{Rook endings!Good placement of the defender} White should bring
his king to the \term{short-side}, with tempo: 67.\K d2! \R h3 68.\K c1 \R h1+ 69.\K d2 \R h3
70.\K c1 \K d5 71.\R g2 \K c4 72.\K b2, and White holds.

@ 8/8/8/8/2k5/1pr5/6R1/3K4 b - - 9 70 | Rd3+?
Black should have brought his king to a3 and then to b2: 70...\K b4! 71.\R g8 \K a3 72.\R a8+
\K b2 73.\K d2 \R c2+ 74.\K e1 \R c5 75.\R b8 \K c2 76.\K e2 b2, and Black wins.

The text move allows 71.\K c1!, and then White holds: 71...\K c3 72.\R g1 \K b4 73.\R g2 \K a3
74.\R g8 \R d4 75.\R a8+ \K b4 (75...\R a4 76.\R xa4+ \K xa4 77.\K b1 draws) 76.\R b8+, and
White gives \term{rear-checks}.\index{Rook endings!Rear checks}

@ 8/8/8/8/2k5/1p1r4/6R1/3K4 w - - 10 71 | Rd2?
71.\K c1! was the only move: the king goes to the short side again.

% Diagram 40 of the 1500-1599 book: Kotas - Sowa
@ 8/8/8/8/5p2/R7/4rk1K/8 w - - 14 65 | Ra4?
White's pieces stand well: the king on the \term{short-side}, the rook on the third rank, where it
keeps the black rook off that rank. 65.\R b3, 65.\R c3 or 65.\R a1 holds the draw; if the pawn
advances, the rook goes back and gives \term{side-checks} along the first two ranks.

65.\R a4? leaves the third rank to attack the pawn, and Black gets the decisive check on it:
65...\K f1+ 66.\K h3 \R e3+! The black rook has \term{cut-off} the white king from its pawn, and
after 67.\K g4 f3 Black wins.\index{Rook endings!Cut-off} Do not improvise the
setup!\index{Rook endings!Good placement of the defender}

@ 8/8/8/8/R4p2/7K/4r3/5k2 b - - 17 66 | f3?
66...\R e3+! would have cut off the white king and won, as shown above. After the text move the
king gets back in time.

@ 8/8/8/8/R7/5p1K/4r3/5k2 w - - 0 67 | Ra3?
The rook move comes too late. Only 67.\K g3! holds: the king comes back next to the pawn, and after
67...f2 68.\K f3 the pawn falls.

% Diagram 42 of the 1500-1599 book
@ 8/8/2K5/2R5/p6r/8/3k4/8 b - - 3 83 | Rh3?
The \term{process-of-elimination} rules this move out,\index{Practical play!Process of elimination}
because White has a simple way to \term{eliminate-the-pawn}:\index{Eliminate the pawn}
83...\R h3 84.\K b5 a3 85.\K b4 \R f3 86.\R a5 \K c2 87.\R xa3, and the game is a draw.

Black wants to push his pawn, and 83...a3 is natural and good.

% Old Diagram 43: Fernandes - Fonseca, too difficult for beginners
@ 8/R7/6k1/2r5/1K6/P7/8/8 b - - 2 50 | skip

% Voges - Chen
@ 6K1/4k1P1/8/8/8/7r/5R2/8 w - - 12 76 | Re2+?!
76.\R e2+ does not throw away the win, but White could have won at once with the standard
\term{lucena}:\index{Rook endings!Lucena's position} 76.\R f4! \R h1 77.\R e4+ \K d6 78.\K f7
\R f1+ 79.\K g6 \R g1+ 80.\K f6 \R f1+ 81.\K g5 \R g1+ 82.\R g4, and we have the position of
\xref{Diagram}{lucena}.

If instead 80...\R g2, then 81.\R e6+ \K c5 82.\R e5+ \K b4 83.\R g5 shows another idea,
\term{bridging}, as in \xref{Diagram}{bridging}.\index{Rook endings!Bridging}

% Diagram 43 of the 1500-1599 book (built 26 Oct): Lakusic - Radosevic, too difficult for beginners
@ 4R3/8/8/8/3r4/4k3/2P5/2K5 b - - 14 74 | skip

% Diagram 47 of the 1500-1599 book: Kaplan - Kara
@ 7k/R4K2/8/5P2/8/8/8/4r3 b - - 4 72 | Rg1?
72...\R g1? is an improvisation.\index{Rook endings!Good placement of the defender} Good is
72...\R b1!: the rook stays on the \term{long-side}, on the first rank. Some sample lines:

73.f6 (73.\R e7 \R a1 74.\K g6 \R g1+ 75.\K f7 \R a1 76.\R e8+ \K h7 77.\R e6 \R a8) 73...\K h7.
Black must never allow White to play \K g6 combined with a mate threat from the a-file, for
example 73...\R b2? 74.\R a5 \R h2 75.\K g6 \R g2+ 76.\R g5, and White wins.

74.\K e7 \R b8. White is threatening f7, and Black must be ready for it. 75.\K f7 (75.f7 \K g7,
and White cannot queen on f8, because the black rook covers the back rank) 75...\R b1 76.\R e7
\R b8, and White cannot make progress.

% Diagram 48 of the 1500-1599 book: Meijerink - Diegel
@ 8/8/8/5pR1/1r3k2/8/5K2/8 w - - 2 66 | Rg3?
66.\R g3? is obviously an oversight. White only has to play by the book, as in
\emph{Good placement}:\index{Rook endings!Good placement of the defender} 66.\R g8!, first the rook
to the eighth rank. 66...\R b2+ 67.\K g1, the king to the \term{short-side}. 67...\K f3 68.\R a8,
and the rook is ready for \term{side-checks} or \term{rear-checks}. 68...\R b1+ 69.\K h2 f4
70.\R a3+, and the side checks begin:\index{Rook endings!Side checks} 70...\K f2 71.\R a2+ \K e3
72.\R a3+ \K e4 73.\R a4+ \K e3 74.\R a3+ \K e4 75.\R a4+, with a draw.

% Diagram 49 of the 1500-1599 book: Danielsson - Haile
@ 8/8/3R4/8/6p1/4r1k1/8/5K2 b - - 7 76 | Rf3+?
White's position is hopeless: his king is on the \term{long-side}, and his rook cannot give
\term{rear-checks} or \term{side-checks}.\index{Rook endings!Good placement of the defender} The text
move helps the white king over to the \term{short-side}: 77.\K g1, and White holds.

The following line shows how to win. Black pushes his pawn slowly but surely: 76...\K h2!
77.\R d2+ \K h1 78.\R d4 \R f3+ 79.\K e2 \R g3 80.\K f2 \K h2 81.\K e2 \R g1 82.\R d7 g3 83.\K f3
\R f1+ 84.\K g4 g2, and the pawn queens.

% Diagram 50 of the 1500-1599 book: Alexandersson - Miranda Herrera.
% The position before 60.Kg1 is the same as the diagram, so the comment on 58.Kg1 is a start note.
@ 8/8/8/6p1/8/1r5k/R4K2/8 w - - 6 58 | start
Black threatens ...\K h2 and then to push the g-pawn. 58.\K g1? does nothing against the threat
58...\R b1+. 58.\R a8! is required (58.\R a7 and 58.\R a6 also hold): the rook goes far away, ready
for \term{side-checks}.\index{Rook endings!Good placement of the defender}

@ 8/8/8/6p1/8/7k/R4K2/1r6 b - - 9 59 | Rb3?
59...g4! is required. Then White's position is hopeless: the king is on the \term{long-side}, and
the rook does nothing. Black can push his pawn slowly, as in \xref{Diagram}{danielsson-haile-76b}.
And 60.\K g1? repeats White's mistake.

@ 8/8/8/6p1/8/1r5k/R7/6K1 b - - 11 60 | g4?
This is wrong: 60...\R b1+! wins, as before. After the text move, 61.\R a1 is a simple way to draw.
Against a g-pawn such a
\vocab[rook endings]{passive-defence}{passive defence}{The defending king stays in front of the pawn
and the rook waits on the first rank. It holds against a knight's pawn, but loses against a bishop's
pawn or a central pawn}\index{Rook endings!Passive defence} is possible:
61.\R a1 g3 62.\K h1 g2+ 63.\K g1 \K g3, and White simply waits with the rook on the first rank,
for example 64.\R d1. Black cannot make progress.

Against an f-pawn or a central pawn, the passive defence loses:
\diagram{8/8/8/8/8/r4pk1/8/1R3K2 b - - 0 1}
1...\R a2 2.\R d1 \R h2 3.\K e1 \R h1+ 4.\K d2 \R xd1+ 5.\K xd1 f2, and the pawn queens.

% Diagram 52 of the 1500-1599 book: Diaz Martin - Sanchez Ramirez
@ 8/4k1K1/6PR/4p3/8/5r2/8/8 w - - 0 57 | Kh7?
White is ready to move his king out of the way and push g7. His candidate moves are 57.\K h7 and
57.\K h8. Alas, he chose the wrong one: after 57.\K h8! \R g3 58.g7! Black cannot stop the pawn.

@ 8/4k2K/6PR/4p3/8/5r2/8/8 b - - 1 57 | Rg3
Now 58.g7 \K f7 59.\K h8 \R xg7 60.\R f6+! is the only move: if Black takes the rook, it is
stalemate.\index{Rook endings!Stalemate} Fortunately White does not need this amusing line to
draw. 58.\R h1 \K f6 59.\R f1+ \K e6 60.g7 e4 61.g8=\Q+ \R xg8 62.\K xg8 \K e5 is also a draw: the
white rook cannot stop the pawn alone, because the white king is too far
away.\index{Rook endings!Rook against pawn}

% Diagram 53 of the 1500-1599 book: Uzun - Kahya
@ 8/8/7r/1P4R1/8/6p1/7k/4K3 w - - 1 73 | Ke2?
Black is threatening ...g2. White has two possible plans. One is to give up the rook for the pawn
and hope that the resulting ending of
\vocab[rook endings]{rook-against-pawn}{rook against pawn}{An ending in which one side has only a
rook and the other only a pawn, usually after a rook was given up for a pawn. The rook wins when
its king is close enough or the defending king is cut off from the pawn; the pawn has drawing
chances only when its own king supports it and the other king is far away} is a
draw.\index{Rook endings!Rook against pawn} As we will see, it is not.

The other plan is to stop the pawn: 73.\K f1! \R f6+ 74.\K e2!, the only move, 74...g2 75.\R h5+
\K g3 76.\R g5+. Because the white king stands on e2, the black king cannot go to f3, and Black
must be satisfied with \term{perpetual-check}; otherwise he loses the g-pawn.

@ 8/8/7r/1P4R1/8/3K4/6pk/8 b - - 1 74 | g1=Q?
74...\R h4! was much stronger. Now the rook has \term{cut-off} the white king on the fourth
rank.\index{Rook endings!Cut-off} Black queens on the next move and forces White to give up his
rook for the pawn, and the ending of \term{rook-against-pawn} is a win for
Black,\index{Rook endings!Rook against pawn} because the white king can never come closer to its
pawn. When the pawn reaches b7, Black plays ...\R h8 and then ...\R b8: Black needs two moves to
attack the pawn, the white king on the third rank needs three to defend it. For example, 75.b6
g1=\Q\ 76.\R xg1 \K xg1 77.b7 \R h8 78.\K c4 \R b8 79.\K c5 \R xb7.

% Diagram 54 of the 1500-1599 book: Angelov - Fil
@ 8/8/8/5p2/8/2R2Pk1/1r6/5K2 w - - 19 59 | Ke1?
There is no point in moving the king to the \term{long-side}.\index{Rook endings!Long side} The
white rook is already on the third rank, so White can set up a
\vocab[rook endings]{philidor}{Philidor position}{The defending king stands in front of the pawn,
and the defending rook on its third rank keeps the attacking king away. To make progress the
attacker has to advance the pawn, and then the rook goes to the far end of the board and checks
from behind}\index{Rook endings!Philidor position} while the black pawn is still on the fifth
rank. To make progress, Black must play ...f3, and that allows White to give
\term{rear-checks}:\index{Rook endings!Rear checks} 59.f4+ \K xf4 60.\R a3 \K g4 61.\K g1 f4
62.\K f1 f3, and now the rook goes back: 63.\R a8 \K g3 64.\R g8+ \K f4 65.\R f8+ \K e3 66.\R e8+
\K f4, with a draw.

@ 8/8/8/5p2/8/2R2Pk1/1r6/4K3 b - - 20 59 | Rb7?
59...f4! is a clear-cut win. It is a
\term{prophylaxis} move:\index{Strategy!Prophylaxis} it stops White's f4+, and Black then wins
the pawn on f3. If White exchanges the rooks, the pawn ending is lost for him, and the rook ending
is probably lost too, since his king is on the long side; the computer confirms it.

We do not need a precise assessment here. The text move gives up the chance to win the pawn, while
59...f4 certainly improves Black's position and gives him real winning chances. So why not take it?

% Diagram 55 of the 1500-1599 book: Wasif - Adhikari
@ 4k3/1K6/1R6/2P3p1/1r6/8/8/8 b - - 2 71 | Rd4?
The black king is on the \term{long-side}, which makes \term{side-checks}
impossible,\index{Rook endings!Long side} and the white rook on b6 already blocks the
\term{rear-checks}. Black's only chance is the second Philidor defence. Compare it with the \term{philidor} of the
previous diagram: according to Dvoretsky, this method was found only in the 20th century, long
after Philidor~\cite{Dvoretsky:2008-endgame-manual}. The black rook goes behind the pawn, and the black king
goes to the back rank, on the same file as the pawn.\index{Rook endings!Second Philidor defence}

71...\R c4! is the only move: 72.\K c6 g4 73.\K d5 \R c1 74.\R g6 \K d8 75.\R xg4 \K c8, and the
setup is complete. 76.\K c6 \K b8 77.\R g8+ \K a7 78.\K d6 \K b7 79.\R g7+ \K c8, and we have the
same setup again. The white king cannot walk away from its pawn, because the rook behind it would
take it.

% Diagram 56 of the 1500-1599 book: Vasiliauskas - Le Boette
@ 8/r7/4Pk2/5P2/4RK2/8/8/8 w - - 9 63 | Rd4?
White has several candidate plans, and we check them one by one.

First, White can push 63.e7, giving up a pawn to exchange the rooks: 63...\R xe7 64.\R xe7 \K xe7.
The pawn ending is a draw, however: 65.\K e5 \K f7! takes the \term{opposition}.

Second, 63.\R d4, the game move, threatens \R d7. Black refutes it with 63...\R a5!: the f5-pawn
falls soon, and the e6-pawn alone cannot promote.

After the \term{process-of-elimination}, we come to \emph{schematic thinking}:\index{Practical play!Schematic thinking}
White brings his king to the queenside and drives the black rook away from the long side. Then he
pushes his pawns. A sample line: 63.\R e1 \R b7 64.\K e4 \R a7 65.\K d5 \K e7 66.\K c6 \R a8
67.\R d1 \R c8+ 68.\K b7 \R f8 69.\R d7+ \K e8 70.\R d5 \R h8 71.f6 \R h1 72.f7+ \K e7 73.\R d7+,
and White wins. The line only shows the idea; you do not need to calculate anything this long.

In the game, Black missed 63...\R a5! and played 63...\K e7?, and White won after all.

% Diagram 57 of the 1500-1599 book: Gloria - Blackman
@ 8/8/8/8/8/r6k/p5p1/R5K1 b - - 1 81 | Kg3?
When you are winning, do not take away the last free squares of the enemy king if your opponent
has no pawn moves left. He may then find a
\vocab{desperado}{desperado}{A piece that is lost anyway, or that has nothing better to do, gives
itself up as dearly as possible: it captures or checks until it is taken. Often taking it is
stalemate, or not taking it allows perpetual check.}, a piece that gives itself up, and save
the game by stalemate,\index{Rook endings!Stalemate}\index{Desperado} as in the game:
82.\R xa2! \R xa2 is stalemate.

In a practical game, you do not need to calculate whether a desperado works when you have another
way to win. Look for that win first. Black could have played 81...\R b3!, threatening \R b1 and
then promoting one of the pawns. This is again a little \emph{schematic thinking}:\index{Practical play!Schematic thinking}
we do not show any moves here, because the idea is enough.

% Diagram 58 of the 1500-1599 book: Ayan - Ustundag
@ 8/8/7R/8/4k1PP/8/1r4K1/8 w - - 5 55 | Kg1?
55.\K g1? is a mistake because White weakens himself: \term{do-not-weaken}. His king is cut off on
the back rank. Again, it is not White's job to find out whether Black has a draw after this move.
He should first check the alternatives that do not weaken his position. 55.\K g3! is good enough:
the king helps the pawns forward with g4-g5, \K g4 and so on.

Black's drawing line is still beautiful: 55...\K f3 56.g5 \R b1+ 57.\K h2 \K g4! 58.\K g2 \R b2+
59.\K f1 \R h2 60.g6. White must push the pawn to make progress; 60.h5 gives up both pawns.
60...\K f3 61.\K g1 \R g2+ 62.\K f1 \R h2 63.\K g1. If the white king goes to the queenside,
the black king simply follows it and creates mating threats. White cannot stop them, because
his rook cannot give checks from there.

63...\R g2+ 64.\K h1 \K f2 65.\R h5 (65.h5 \R g4 66.\K h2 \R h4 mate) 65...\R g1+ 66.\K h2 \R g2+
67.\K h3 \R g3+, and Black gives \term{perpetual-check}.\index{Perpetual check}

@ 8/8/7R/6P1/7P/5k2/7K/1r6 b - - 2 57 | Rb2+?
Now the white king escapes again. 57...\K g4! was the only move to draw.

% Diagram 59 of the 1500-1599 book: Kahya - Incecik
@ 5R2/8/8/8/5kp1/1r5p/7K/8 b - - 5 72 | Ke3?
This move allows White to defend. An easier win was to retreat the king, then bring the rook to
g2 with check and push h2: 72...\K g5 73.\R g8+ \K f5 74.\R f8+ \K e6 75.\R f1 \R b2+ 76.\K g3
\R g2+ 77.\K f4 h2, and Black wins.

@ 6R1/8/8/8/4r1p1/6Kp/8/4k3 w - - 14 77 | Rf8?
White has several moves here, but this is not one of them. Cutting off the black king does not
defend, because it does not stop Black from playing \R e2, \R g2+ and then h2, as in the last
comment. Checks from the side guarantee a draw: 77.\R b8 \K f1 78.\R b1+ \R e1 79.\R b2 \R e3+
80.\K xg4 \K g1 81.\R b1+ \K f2 82.\R b2+, and the checks go on.\index{Rook endings!Side checks}

% Diagram 60 of the 1500-1599 book: Androni - Markaki
@ 5R2/6rp/4k3/6PP/5K2/8/8/8 b - - 0 72 | Ke7?
72...\R f7+! 73.\R xf7 \K xf7 is an immediate draw. As we know, the g- and h-pawns against an
h-pawn on its original square are a draw: the defender only needs to move his king between g8 and h8, as in
\xref{Diagram}{ince-tokat-87b}.\index{King and pawn endings!G- and h-pawns against an unmoved h-pawn}

Black probably went wrong because of a common guideline: when you are better, exchange pieces
and keep the pawns; when you defend, exchange pawns and keep the pieces. We do not want to spend
energy on whether the guideline is misleading. In chess, when you find a concrete line that
reaches your goal, you can break every guideline you have heard.\index{Practical play!Concrete line beats guidelines}

@ 5k2/R5rp/8/4K1PP/8/8/8/8 w - - 7 76 | Rxg7?
White makes the same mistake. If the guideline is not misleading, why do two players make the
same mistake? 76.\R a8+! wins: 76...\K f7 77.\K f5 \K e7 78.g6 hxg6+ (78...h6 79.\K e5 \K d7
80.\R a7+ wins the rook) 79.hxg6 \R xg6 80.\K xg6. The resulting rook ending is hopeless for
Black, because none of his pieces is well placed: his king is on the long side, and his rook is
passive.

% Diagram 61 of the 1500-1599 book: Majstrova - Petruzalek
@ 8/5k2/8/2p5/1p6/2r3K1/P3R3/8 w - - 2 63 | Kf2?
One important principle in the endgame is \vocab[strategy]{king-activity}{king activity}{In the
endgame the king is a strong piece. Bring it to the centre, or to where the play is, before the
opponent's king gets there.}.\index{Strategy!King activity}
On f2 the king defends nothing on the back rank. 63.\K f4 is certainly better, because the king
is more active.

@ 8/5k2/8/2p5/1p6/2r5/P3RK2/8 b - - 3 63 | c4?
Do not rush! (\vocab[strategy]{do-not-rush}{do not rush}{When the opponent has no real counterplay, there is no
hurry. First take away his resources and improve your pieces, then act.})\index{Strategy!Do not rush} Black should improve his position first, before
taking decisive action: 63...\K f6! 64.\K e1 \R d3 65.\R e8 \R d5 66.\R h8 \K e5, and Black's
position is dominating.

@ 8/5k2/8/8/1pp5/2r5/P3RK2/8 w - - 0 64 | Ke1?
White misses his chance: 64.\R b2! takes the pawn and draws at once.

@ 8/5k2/8/8/1pp5/2r5/P3R3/4K3 b - - 1 64 | b3?
Again, \termas{do-not-rush}{do not rush}!\index{Strategy!Do not rush} 64...\R d3! wins.

% Diagram 62 of the 1500-1599 book: Diaz Martin - Sanchez Ramirez
@ 6R1/8/4k1K1/4p3/8/2r2PP1/8/8 w - - 3 47 | Re8+?!
According to the computer, 47.\R e8+ throws away the win. We do not want to be too harsh on this
move, because it is a good practical move: Black faces a choice, and he goes wrong.

A simpler plan is 47.\R f8!. White's only weakness, the f3-pawn, is protected, the black king is
\term{cut-off},\index{Rook endings!Cut-off} and White can push his g-pawn with the help of his king.
Black's only try to spoil this plan is 47...e4, but after 48.\R e8+ \K d6 49.\R xe4 \R xf3 50.g4
White wins.

In a complex ending, it is important to evaluate the position a few moves ahead. Here we reach a
rook and g-pawn against a rook. The black king is cut off far away, and the white rook on the fourth
rank is ready for the \term{lucena}, where it will block Black's \term{rear-checks}.\index{Rook endings!Lucena's position}

@ 4R3/8/4k1K1/4p3/8/2r2PP1/8/8 b - - 4 47 | Kd5?
Black faces a choice and goes wrong. 47...\K d6! draws: 48.\R f8 e4!. Again, Black must create
counterplay before White pushes his g-pawn. 49.\R f6+ (49.fxe4 does not win either, and the
\term{process-of-elimination} rules it out: 49...\R xg3+ 50.\K f6 \R f3+ 51.\K g7 \R xf8 52.\K xf8
\K e5)\index{Practical play!Process of elimination} 49...\K e5 50.f4+ \K d5 51.g4, and Black has
his own passed pawn. It is important to create counterplay like this, even when you are a pawn
down.\index{Practical play!Counterplay}

After 47...\K d5? White wins: 48.\R f8 e4 49.fxe4+ \K xe4 50.g4.

% Diagram 64 of the 1500-1599 book: Rxc5? and a skewer
@ 6R1/8/8/2P2k2/2r2p1P/4pK2/8/8 b - - 2 53 | Rxc5?
Black makes the only move that loses: after 54.\R g5+ he loses his rook. Such a simple error can
be avoided by the \term{process-of-elimination}.\index{Practical play!Process of elimination}

% Diagram 65 of the 1500-1599 book: Majstrova - Petruzalek
@ 8/5R2/1p6/2p1kP2/2r5/6K1/P7/8 b - - 2 57 | b5?
This misses the point: Black can take the a2-pawn. 57...\R a4! 58.f6 \R xa2 59.\R f8 \K e6
60.\K g4 \R d2 61.\K g5 \R g2+ 62.\K f4 \R f2+ 63.\K g5 c4 64.\R c8 \R f5+ 65.\K g4 \R c5 wins.

@ 5R2/4kP2/8/2p5/1pr5/6K1/P7/8 w - - 1 61 | Re8+?
61.\R a8! draws: 61...\K xf7 62.\R a6!. White cuts off the black king.\index{Rook endings!Cut-off}
62...\K e7 63.\K f3 \K d7 64.\K e3 \K c7 65.\K d3 \R c1 66.\R g6 c4+ 67.\K d2 \R a1 68.\R g4.
White goes after the pawns. 68...\R xa2+ 69.\K c1 \R a6 70.\R xc4+ \R c6 71.\R xc6+ \K xc6
72.\K c2 \K b5 73.\K b1 b3 74.\K b2 \K b4 75.\K b1 \K c3 76.\K c1 b2+ 77.\K b1 \K b3, and it is
stalemate.

@ 8/5k2/8/2p5/1pr5/6K1/P3R3/8 b - - 1 62 | Rc3+?
62...\K f6! keeps the king active.\index{Strategy!King activity}

% Diagram 66 of the 1500-1599 book
@ 8/5p2/6k1/4K1Pp/4P3/5r2/8/2R5 w - - 2 71 | Rg1?
There is no point in protecting the g5-pawn: it is lost anyway. The goal is to save the half point
by looking for a clear-cut draw, not to protect material.\index{Practical play!Clear-cut draw}
71.\R h1! \K xg5 72.\R g1+ \K h6 73.\R h1 is a simple draw, because Black cannot push his h-pawn.

% Diagram 67 of the 1500-1599 book
@ 8/8/5R1P/3k2Pr/4p1K1/5p2/8/8 b - - 1 73 | Rh2?
We know that a rook needs room to give checks. That is why the rook should stay on the
\term{long-side}.\index{Rook endings!Long side} 73...\R h1! draws: 74.g6 \R g1+ 75.\K h5 \R h1+
76.\K g5 \R g1+, and Black gives \term{perpetual-check}.\index{Perpetual check}

@ 8/8/5R1P/3k2P1/4p1K1/5p2/7r/8 w - - 2 74 | Kf4?
Pushing the pawn wins at once: 74.g6! \R xh6 75.g7 \R xf6 76.g8=\Q+.

@ 8/8/5R1P/3k2P1/4pK2/5p2/7r/8 b - - 3 74 | Kd4?
Perpetual check draws at once: 74...\R h4+ 75.\K e3 \R h2 76.\K f4 \R h4+.\index{Perpetual check}

@ 8/8/3R3P/6P1/4pK2/2k2p2/7r/8 w - - 6 76 | Rd1?
76.\K xe4! wins: 76...f2 77.\R d1 \K c2 78.\R f1 \K c3 79.\K e3 \R h5 80.\K f4 \R h2 81.\K f3
\R h3+ 82.\K xf2.

@ 8/8/7P/6P1/4pK2/2k2p2/7r/3R4 b - - 7 76 | Rf2?
This move has no target at all. 76...\K c2! helps the pawn to promote and draws: 77.\R f1 \K d3
78.\K g3 \R g2+ 79.\K f4 f2 80.h7 \R h2 81.g6 e3 82.g7 e2 83.h8=\Q{} \R xh8 84.gxh8=\Q{} exf1=\Q+.
The computer says draw, but Black has an extra pawn, and in a practical game the opponent can
go wrong under time pressure or fatigue. Black even has some chances to win here, certainly more
than in the game.

% Diagram 59, continued: the defence against connected knight's and rook's pawns
@ 5R2/8/8/8/6p1/1r2k2p/7K/8 w - - 6 73 | Kg3!
A typical defence against connected knight's and rook's pawns: White puts his king on g3, in
front of the pawns. He threatens to take the g4-pawn, or to set up a fortress along the second
rank. Black cannot break through.\index{Rook endings!Connected knight's and rook's pawns}

@ 5R2/8/8/8/1r4p1/4k1Kp/8/8 w - - 8 74 | Re8+
The easiest draw is 74.\R a8 and then \R a2. Black cannot make progress against this
\vocab[strategy]{fortress}{fortress}{A defensive set-up that the stronger side
cannot break, even with extra material. The defender only has to keep it.},
for example 74...\K d2 75.\R a2+ \K d3 76.\R a3+ \K c2 77.\R a2+ \K b3 78.\R f2. The black king
can never stay on the second rank.\index{Fortress}\index{Rook endings!Connected knight's and rook's pawns}

% Diagram 70 of the 1500-1599 book
@ 8/8/5Pk1/7R/r5PP/8/8/5K2 b - - 0 56 | Rf4+?
Black is three pawns down. 56...\R xg4? 57.\R g5+ \R xg5 58.hxg5 is hopeless. His one hope is
rook against rook and connected knight's and rook's pawns, as in
\xref{Diagram}{kahya-incecik-72b}.\index{Rook endings!Connected knight's and rook's pawns} For that,
his king must stand in front of the pawns. In the game, his king is \term{cut-off}
instead.\index{Rook endings!Cut-off}

56...\K xf6! draws: 57.g5+ (57.\R g5 \R a1+ 58.\K e2 \R a2+ 59.\K d1 \R a1+ 60.\K c2 \R a4
61.\R g8 \K f7 62.\R g5 \K f6. White can never take his rook off the g-file without losing
material. If he does, the ending is a draw, because his king is far away and his last pawn
falls.) 57...\K g7 58.\R h6 \R a2 59.h5 \R a5 60.\R g6+ \K h7 61.\K g2 \R b5.

This is a typical defence. Black keeps the g5-pawn under attack, so the white rook cannot leave
it. As long as the white king stays behind the fourth rank, Black just waits. When the king
reaches the fourth rank, Black gives checks. If White tries to escape the checks, Black exchanges
the rooks, and the pawn ending is a draw. For example, 62.\K g3 \R a5 63.\K f4 \R a4+ 64.\K e5
\R a5+ 65.\K d6 \R a6+ 66.\K e5 \R xg6 67.hxg6+ \K xg6 68.\K f4.

% Diagram 63 of the 1500-1599 book: Yassibag - Tahaoglu
@ 8/8/8/1pR5/p2K4/P7/1r1k4/8 w - - 0 80 | Ke4?
We can use the \term{process-of-elimination} here.\index{Practical play!Process of elimination}
Black will play b4, and after the pawns are exchanged White must defend rook against rook and
rook's pawn. It helps to keep the black king \term{cut-off},\index{Rook endings!Cut-off} so that it
cannot help its pawn, and the white king should stay close to the pawn. So 80.\K d5! is the right
move. A possible line: 80...b4 81.axb4 \K d3 82.\R c1 \R xb4 83.\K c5 \R b3 84.\R h1 \K c2 85.\K c4
\K b2 86.\R h2+ \K a1 87.\R h1+ \K a2 88.\R h2+ \R b2 89.\R h1 \R b8 90.\R h2+. White again uses
\term{side-checks};\index{Rook endings!Side checks} his king stays on c4 so that it does not block
the checks.

You do not need to calculate this whole line. We used principles to find the right move.

% Diagram 68 of the 1500-1599 book, left out
@ 8/1k6/7R/1P6/8/1P6/1r5P/6K1 w - - 1 71 | skip

% Diagram 69 of the 1500-1599 book
@ k7/8/P7/1K3r2/4RP2/P7/8/8 w - - 9 67 | Re5?
White has a few candidate plans. One is the game move: giving up the f4-pawn. The
\term{process-of-elimination} should have ruled it out,\index{Practical play!Process of elimination}
because the black king is already on a8, and the doubled a-pawns cannot win. One simple defence
is to keep the rook on the back rank and the king on a8. Once White plays a7, the black king
cannot move, and Black can give checks with his rook as a \term{desperado}.\index{Desperado}

The other plan is to keep the f4-pawn. Black's defence against the f-pawn alone is hopeless,
because his king is on the \term{long-side}.\index{Rook endings!Long side} White can play \K b4 and
\K b3, then \R b4, cutting off the black king,\index{Rook endings!Cut-off} and slowly bring his king
to the king side to escort the pawn. The black rook alone cannot stop the pawn without giving
itself up for it, and then White's remaining rook decides the game.

% Diagram 71 of the 1500-1599 book
@ 7k/7p/R7/4K3/3p4/p4r2/8/8 b - - 3 63 | d3?
Black is three pawns up, and the a3-pawn is his trump. Nothing forces him to give it up and leave
his d-pawn vulnerable, as in the game. With such an advantage it is practical not to solve the
problem tactically: the risk of a miscalculation is not worth it. A more prosaic way is better:
63...h5! 64.\R h6+ (64.\K e4 \R h3 65.\K xd4 h4 66.\R h6+ \K g7 67.\R a6 \R g3) 64...\K g7
65.\R xh5 a2 66.\R h1 \R a3 67.\R a1 d3, and Black wins.

% Diagram 72 of the 1500-1599 book: Mourogiannis - Klabatseas
@ 8/1R6/1p6/6kP/P7/6P1/1r6/6K1 w - - 4 73 | a5?
Connected passed pawns are a real asset, and there is no point in giving them up. 73.g4! wins.

@ 8/1R6/1P6/7k/8/6P1/1r6/6K1 b - - 0 74 | Kg6?
Passive defence cannot hold here.\index{Rook endings!Passive defence} White's plan is simple: his
king walks via a6 to b7 and frees his rook, and then the b-pawn cannot be stopped. Black must play
actively: 74...\K g4! draws.

@ 1R6/6k1/1P6/8/8/6P1/1r6/6K1 w - - 3 76 | b7?
This is an instructive mistake. With the pawn on b7, the white king can no longer hide on b7.
Black's defence is now very simple: his king stays on g7 (or h7). For example, 76...\K f7? loses to
77.\R h8 \R xb7 78.\R h7+, winning the rook. 76.\K f1! wins.

When the white king walks to the queenside, it cannot escape the checks.

@ 1R6/1P4k1/1r4P1/8/8/4K3/8/8 b - - 6 83 | end
The game is objectively still a draw, but somehow White won it. The draw is elementary:
83...\R b3+ 84.\K d4 \R b1 85.\K c5 \R c1+ 86.\K b4 \R b1+ 87.\K a4, and the checks go on. Black
must not take the g6-pawn, though: after ...\K xg6 White wins with \R g8+ and b8=\Q.

% Diagram 73 of the 1500-1599 book, left out
@ 8/5r2/1R6/3P2k1/2P5/2K4p/P7/8 w - - 0 47 | skip

% Diagram 74 of the 1500-1599 book
@ 5k2/5p2/R4p2/1pK2P2/4r3/8/8/8 w - - 0 59 | Rxf6?
White must make a decision here. He can take the f6-pawn, as in the game, and then defend
against two passed pawns. Or he can take the more dangerous b5-pawn and defend against the
doubled f-pawns. Whether the position against the doubled f-pawns can be held is secondary:
it is certainly much better than the other one, and we have tried our best. In fact, the
computer says that 59.\K xb5! draws. After the text move, the position is hopeless.

% Diagram 74 of the 1500-1599 book (after Diagram 73 was left out)
@ 8/4k1p1/6P1/4p3/4rp2/3K4/5R2/8 b - - 3 51 | Kf6?
Black forgets that his rook is hanging, and he loses a winning position.

% Diagram 75, left out
@ 8/7p/5k2/5P2/4K3/2R5/2Pr3P/8 w - - 3 47 | skip

% Diagram 76 of the 1500-1599 book
@ 4k3/1R6/2P5/p1K4p/1p6/8/3r4/8 w - - 0 56 | Rh7?
Now the connected passed pawns become very dangerous. 56.\R b5! is a typical defence against
connected passed pawns: neither pawn can move. 56...h4 57.\R xa5 b3 58.\R b5 b2 59.\R b8+ \K e7
60.c7 \K f6, and White holds.

@ 4k3/2P4R/3K4/p6p/8/1p6/2r5/8 b - - 2 58 | b2?
\termas{do-not-rush}{Do not rush!}\index{Strategy!Do not rush} Black misses that White can bring his rook to b8,
take on b2 and exchange the pawns. The ending after that is in fact a draw, but that is
secondary: why allow the opponent such a chance? \term{prophylaxis} would be much wiser:\index{Strategy!Prophylaxis}
58...a4! keeps Black's winning chances.

@ 1R6/2P2k2/3K4/p6p/8/8/1pr5/8 b - - 3 60 | a4?
Now Black loses. 60...\R d2+! 61.\K e5 \R e2+ 62.\K f4 (62.\K d4 \R c2) 62...\R c2 63.\R xb2
\R xc7 64.\R b5 \R c4+ 65.\K e3 h4 66.\R xa5, and White holds. The text move should have been
ruled out by the \term{process-of-elimination}.\index{Practical play!Process of elimination}

% Diagram 76 of the 1500-1599 book
@ k3r3/P1R5/1K4p1/5PP1/8/8/8/8 w - - 0 76 | fxg6?
An oversight. The text move allows Black to exchange the pawns.\index{Defensive resources!Exchanging or removing the pawns}
76.f6! wins, because Black cannot stop \R e7, f7 and \R e8.

% Diagram 77 of the 1500-1599 book
@ 8/8/7p/6k1/P2R2P1/r4K1P/8/8 w - - 4 49 | Kg2?
The king must become active, even at the cost of material.\index{Strategy!King activity}
49.\K e4! \R xh3 50.\K d5 \R a3 51.\K c6 \R c3+ 52.\K b5 \R b3+ 53.\R b4 wins.

@ 8/8/7p/6k1/P2R2P1/r6P/6K1/8 b - - 5 49 | Ra2+?
The black rook cuts off the white king, and there is no point in releasing it with an
\vocab[practical play]{unnecessary-check}{unnecessary check}{A check that achieves nothing.
It often drives the enemy king to a better square, or costs a tempo.}.\index{Typical errors!Unnecessary check}
49...\K h4! 50.\R d5 \R g3+ 51.\K f2 \K xh3 52.\R d4 \R a3 leads to an ending that is indeed a
draw. The text move gives White an easy win.

% Diagram 44 of the 1500-1599 book (Fernandes - Fonseca), left out
@ K7/P3k3/8/8/1r6/8/2R5/8 w - - 3 64 | skip

% Diagram 46 of the 1500-1599 book (Hadzala - Deng), left out
@ 8/5k2/8/7r/P1KR4/8/8/8 w - - 3 60 | skip

% Diagram 45 of the 1500-1599 book: Hadzala - Deng
@ 8/5k2/8/3R3r/P1K5/8/8/8 b - - 4 60 | Rh6?
Against a rook's pawn it can be tricky, because the defender's king is always on the
\term{long-side}.\index{Rook endings!Long side} Black should have seized the moment to get a
\vocab[practical play]{clear-cut-draw}{clear-cut draw}{A line that reaches a plainly drawn position
by force, so that nothing is left to calculate or to go wrong.}\index{Practical play!Clear-cut draw}
by giving \term{rear-checks}:\index{Rook endings!Rear checks} 60...\R h4+ 61.\K b5 \K e6 62.\R d3
\R h1 63.a5 (63.\K c4 \R a1 64.\K b3 \R b1+ 65.\K a3 \R a1+ 66.\K b3 \R b1+ 67.\K a2 \R b8: White
cannot get rid of the \term{perpetual-check})\index{Perpetual check} 63...\R b1+ 64.\K a4 \R a1+
65.\K b3.

@ 8/5k2/7r/3R4/P1K5/8/8/8 w - - 5 61 | Kb5?
This move allows \term{rear-checks} again.\index{Rook endings!Rear checks} 61.a5 wins, because
White can use \term{bridging}\index{Rook endings!Bridging} to block the rear checks: 61...\R h1
62.a6 \R a1 63.\K b5 \R b1+ 64.\K c6 \R c1+ 65.\K b6 \R b1+ 66.\R b5.

@ 8/5k2/7r/1K1R4/P7/8/8/8 b - - 6 61 | Ke7?
Black misses an opportunity: 61...\K e6 62.\R d1 \R h2, the same defence again. Black prepares
\term{rear-checks}:\index{Rook endings!Rear checks} 63.a5 \R b2+ 64.\K c6 \R c2+ 65.\K b5 \R b2+
66.\K c4 \R a2.

% Diagram 51 of the 1500-1599 book: Faulhaber - Sreyas
@ 2r5/4R3/5k2/8/2P5/2K5/8/8 w - - 14 65 | Rd7?
White could win by \termas{cut-off}{cutting off} the black king:\index{Rook endings!Cut-off} 65.\R e1 \R a8 66.\K b4
\R b8+. Black must try \vocab[rook endings]{frontal-checks}{frontal checks}{The defending rook
checks the attacking king from the front, so that the pawn cannot move forward. Otherwise the
defence is lost.}.\index{Rook endings!Frontal checks} When White's pawn moves forward, Black can no
longer defend, because his king is cut off on the \term{long-side}.\index{Rook endings!Long side}
67.\K c5 \R c8+ 68.\K b5 \R b8+ 69.\K a6 \R c8 70.\R e4 \K f5 71.\R h4 \K f6 72.\K b7 \R c5
73.\K b6 \R c8 74.c5, and the win should be easy.

% Diagram 196 of the 1500-1599 book: De Paula - Miguel
@ 8/1R3bk1/8/6K1/8/8/8/8 b - - 4 56 | Kf8?
Black should have used the \term{process-of-elimination}\index{Practical play!Process of elimination}
to exclude the text move. 56...\K g8 is the only move left, because the king must protect the
bishop. After the text move, 57.\K f6 threatens a
\vocab{back-rank-mate}{back-rank mate}{Mate on the edge of the board: the rook checks
along the last rank, and the king cannot leave it, because the enemy king or its own pieces take
away the squares in front of it.}.\index{Rook!Back-rank mate} 57...\B e8 58.\R b8, and Black
is in \term{zugzwang}:\index{Zugzwang} he must lose his bishop and will soon be mated.

% Diagram 197 of the 1500-1599 book: Szucs - Borhy
@ 8/8/8/4p3/RP6/2b2K2/2k5/8 b - - 0 63 | Kd3?
White's rook and pawn are clumsy together. Black should have won a pawn with 63...\K b3.

% Diagram 198 of the 1500-1599 book
@ 8/8/KP6/8/1r3k2/8/7p/7B b - - 0 57 | Ra4+?
White threatens \K a7, b7 and b8, giving up the pawn for the rook; the remaining ending is a draw.
Black must react immediately: 57...\K e5 58.\K a7 \K d6 59.b7 \K c7, threatening a
\term{back-rank-mate}.\index{Rook!Back-rank mate} White cannot promote. 60.\K a6 \R b6+ 61.\K a7
\R b1, another back-rank mate threat. 62.b8=\Q+ \R xb8.

% Diagram 199 of the 1500-1599 book: Kwiek - Kozlinski
@ 5k2/2R5/5K2/7p/6bP/8/8/8 w - - 5 62 | Rb7?
The win is less obvious. To win the game, White must try to win the h5-pawn, and this is possible
thanks to \term{zugzwang}.\index{Zugzwang} The remaining ending is winning, because White's
h-pawn has not crossed the middle of the board yet (see \cite{Dvoretsky:2008-endgame-manual} for
more). 62.\R e7 \B h3 63.\R e1 \B d7 64.\K g5 \B e8 (64...\B g4 65.\K g6: Black is in zugzwang and
must lose a pawn) 65.\R e3 \B a4. Black must abandon the h5-pawn (65...\B f7 66.\K f6 \K g8
67.\R g3+ \K f8 68.\R a3 \K e8 69.\R a8+ \K d7 70.\K xf7). 66.\K xh5 \B d1+ 67.\K g5. Now White has
won the pawn.

His plan: Black's king must be \termas{cut-off}{cut off}\index{Rook endings!Cut-off} from the
g- and h-files. White places his rook on g5 and plans the manoeuvre \K h5--g4. Only afterwards does
he push h5. Note that if White plays h5 too early, the pawn on h5 blocks the king's
route~\cite{Dvoretsky:2008-endgame-manual}. 67...\B c2 68.\K h6 \K f7 69.\R g3 \B b1 70.\R g7+
\K f8 71.\R g5 \B a2 72.\K h5 \B b1 73.\K g4 \B e4 74.h5 \K f7 75.h6 \B b1 76.\R g7+ \K f6
77.\K f4 \B d3 78.\R c7 \K e6 79.h7.

@ 4k3/8/5K2/7p/5R1P/5b2/8/8 b - - 68 93 | Bd1?
Black misses that White can threaten a \term{back-rank-mate}\index{Rook!Back-rank mate} and
win the bishop. 93...\B g4 still holds.

% Diagram 200 of the 1500-1599 book: De Leon - Manteiga
@ 2r5/4k2K/7B/6PP/8/8/8/8 b - - 2 66 | Rc4?
To win the game, White must move the king and the bishop away and then push the pawns. Black first
forces the g-pawn to move forward, and then blocks it with king and rook: 66...\R c5 67.g6 \R c8
68.\B e3 \K f6 69.h6 \R d8 70.\B c1 \R b8, the simplest move, allowing White no check on b2
(70...\R a8 also draws: 71.\B b2+ \K e6, because Black can give \term{perpetual-check};\index{Perpetual check}
but 70...\R e8? 71.\B b2+ \K e6 72.\K g7, and Black has no more checks). We may remember the
\term{three-file-rule}\index{Rook!Three-file rule} from the rook endings: the rook must be three
files or ranks away from the king to give perpetual check.

@ 8/4k2K/7B/6PP/2r5/8/8/8 w - - 3 67 | g6?
White should have won by moving the bishop away: 67.\B g7 \R c6 68.\B d4 \R c8 69.\K g6. Black
cannot stop the h-pawn, and the remaining g-pawn wins the pawn ending.

@ 8/4k2K/6PB/7P/2r5/8/8/8 b - - 0 67 | Rc7?
It is not too late to play 67...\R c8, and we are back in the line that holds.

% Diagram 201 of the 1500-1599 book: Alexakis - Stefanatos
@ 8/8/1k1p4/8/7p/2R3bP/4K3/8 b - - 2 52 | Kb7?
Black must reckon that White can play \R xg3. There is no pawn race here, since White can collect
the g-pawn on the way and defend against the d-pawn. Black's king is two moves outside the
\termas{rule-of-the-square}{square},\index{King and pawn endings!Rule of the square} so he cannot stop
White from promoting. \termas{prophylaxis}{Prophylaxis}\index{Strategy!Prophylaxis} such as 52...\B f4 holds the
game. In the game, both sides miss the move.

@ 8/1k6/3p4/8/7p/2R3bP/4K3/8 w - - 3 53 | Kf3?
53.\R xg3! hxg3 54.h4 wins: the pawn runs first, and the king collects the g-pawn later.

% Diagram 202 of the 1500-1599 book: Batzoli - Louka
@ 8/8/8/3r4/1P4p1/P3B1k1/4K3/8 b - - 0 68 | Kg2?
\termas{do-not-lose-tempo}{Do not lose a tempo in a pawn race}.\index{Pawn races!Do not lose a tempo}
68...\K h3 69.\B f4 (69.\B g1 g3 70.a4 \K g2 71.\B b6 \K h1 72.\K f3 g2; 69.a4 \R e5 70.a5 g3 71.a6 g2
72.\K d3 \R xe3+ 73.\K xe3 g1=\Q+) 69...g3 70.\B b8 g2 71.\B a7 \R d8 72.\B c5 \K h2 73.\K f3 \R d3+
74.\K e2 \R xa3.

@ 8/8/8/3r4/PP6/4B1p1/4K1k1/8 w - - 0 70 | b5?
This is an instructive mistake. Which pawn should White push? It turns out that White should push
the a-pawn first, because it can go to a7 supported by the bishop. If he pushes the b-pawn first,
as in the game, he must then push the a-pawn to support it, again losing a tempo in the pawn race:
70.a5 \R e5 71.a6 \K h2 72.a7 \R xe3+ 73.\K xe3 g2 74.a8=\Q{} g1=\Q+.

% Diagram 203 of the 1500-1599 book: Headlong - Denning
@ 8/6p1/3k4/5P2/2B1PK1r/8/8/8 w - - 1 56 | Ke3?
White must reckon that Black can give up the rook for a pawn and the bishop, \termas{transposing}{transposing
into a pawn ending}\index{Transposing into a pawn ending} that wins for him. So White must
\termas{king-activity}{activate his king}.\index{Strategy!King activity} The text move allows
...\K e5 and ...\R h3+, and the white king must retreat to more passive squares. 56.\K f3 \K e5
57.\B d5 \K f6 58.\B e6 (58.\K g3 \K g5 59.e5 \R d4 60.\B f7 \K xf5; 58.\B b7 \K g5 59.\K e3 \R h6
60.\B c8 \R h3+ 61.\K d4 \R h8 62.\B b7 \R b8 63.\B d5 \K f4; 58.\B c6 \K g5 59.\K e3 \R h6 60.\B d7
\R h3+ 61.\K d4 \R h8) 58...\K g5 59.e5 \R f4+ 60.\K e3.

@ 8/6p1/8/3BkP2/4P3/7r/3K4/8 b - - 10 60 | Rh1?
60...\K d4 61.\K d1 (61.\K e2 \R e3+) 61...\R e3 62.f6.

% Diagram 204 of the 1500-1599 book: Szucs - Borhy
@ 8/8/8/4p3/R7/1Pb3K1/2Pk4/8 w - - 3 62 | Kf3?
White must reckon that his connected pawns are weaknesses, and he must get rid of them fast. The
position is also a pawn race, and there is no time to lose.\index{Pawn races!Do not lose a tempo}
62.\R g4 \K xc2 63.b4 \B d4 64.b5 \K d3 65.\R xd4+: White wins by \termas{transposing}{transposing
into a winning pawn ending}.\index{Transposing into a pawn ending}

% Diagram 205 of the 1500-1599 book: Duru - Atabey
@ 5r2/8/1p1k2KB/2p2P2/8/8/8/8 b - - 5 54 | Re8?
This is another pawn race, and there is no time to lose.\index{Pawn races!Do not lose a tempo}
Black's rook simply waits on the back rank for the f-pawn, and it gives itself up for the pawn
when it promotes. The remaining ending of connected pawns against bishop should be winning. The
text move simply wastes a tempo (\term{do-not-lose-tempo}), because the rook is attacked on e8
when the pawn reaches f7. In a practical game one does not need to calculate whether the waiting
move still wins. We show the variations of the ending of bishop against connected pawns:
54...\R a8 55.f6 b5 56.\K f5 \K d5 57.f7 b4 58.f8=\Q{} \R xf8+ 59.\B xf8 \K d4 60.\B g7+ \K d3
61.\B f8 b3 62.\B d6 c4 63.\B e5 c3 64.\B g3 b2.

@ 4r3/8/1p1k2KB/2p2P2/8/8/8/8 w - - 6 55 | Bf4+?
55.f6 b5 56.f7 \R a8 57.f8=\Q+ \R xf8 58.\B xf8+ \K d5 59.\K f5 \K d4 60.\K e6 b4 61.\B g7+ \K d3
62.\K d5 c4 63.\K c5 b3 64.\K b4.

% Georgiev - Simeonov
@ 8/4k3/6pp/8/r7/3K1BP1/7P/8 b - - 3 49 | Ra2?
White should not be allowed to play h4, because afterwards Black cannot make further progress.
49...g5 is \term{prophylaxis}.\index{Strategy!Prophylaxis} White's weakness on h2 forces him to play
h3. The remaining ending of rook and rook's pawn against bishop and rook's pawn is favourable for
Black, since White's king is in the
\vocab[rook vs bishop]{dangerous-corner}{dangerous corner}{In rook against bishop, the corner of the
same colour as the bishop. There the defending king can be mated; the safe corner is the one of
the other colour.}.\index{Rook vs bishop!Dangerous corner} The line is quite long, but the idea is
simple: Black drives White's king away with \term{zugzwang}\index{Zugzwang} and the threat of a
\term{back-rank-mate}.\index{Rook!Back-rank mate} When White's king is \termas{cut-off}{cut off} on the
c-file,\index{Rook endings!Cut-off} Black comes back to the pawn and gives up the rook for the bishop.
See \cite{Dvoretsky:2008-endgame-manual} for more.

49...g5 50.\K e3 \R a2 51.\B e2 \K f6 52.h3 \K e5 53.\B h5 \R h2 54.\B g4 h5 55.\B d7 \R g2 56.\K f3
\R c2 57.\B e8 \R c3+ 58.\K f2 h4 59.gxh4 gxh4 60.\B d7 \K f4 61.\B g4 \R c2+ 62.\K e1 \K g3
63.\B f5 \R b2 64.\B d7 \K f3 65.\K d1 \K e3 66.\B g4 \K f2 67.\B f5 \R b5 68.\B e6 \R b6 69.\B f5
\R d6+ 70.\K c1 \K g2 71.\K c2 \R f6 72.\B g4 \R g6 73.\B f5 \R g3 74.\B e6 \R xh3 75.\B xh3+ \K xh3
76.\K d2 \K g2.

% Diagram 208 of the 1500-1599 book: Chong Chor Yuen - Lam
@ 4R3/8/8/5p2/2pb1k2/5P2/6KP/8 w - - 0 51 | Re2?
Too passive. The rook must be activated:
\vocab[strategy]{active-pieces}{active pieces}{Pieces that attack, control many squares and support
their own pawns, instead of only defending. A passive piece often becomes a weakness itself.}.\index{Strategy!Active pieces}
51.h4 would be the first idea, but 51...\B f6! frustrates it: 52.h5 c3 53.h6 c2 54.\R c8 \K g5
55.h7 \K g6. Since we know that ...\B f6 frustrates our plan, we can play 51.\R e6, which is
\term{prophylaxis}.\index{Strategy!Prophylaxis} 51...\K g5 (51...c3 52.h4 \B e3 53.\R c6. Pay
attention to how active the rook is: on one hand it stops Black's pawn from moving forward, on the
other it also supports its own pawn. 53...\B d4 54.h5 \K e3 55.h6 \K d2 56.\R d6 \K d3 57.h7 c2
58.h8=\Q{} c1=\Q{} 59.\Q xd4+ \K c2 60.\Q d3+ \K b2) 52.\K g3 f4+ 53.\K g2 \B a1 54.\R c6 c3
55.\R c5+.

@ 8/8/8/5pk1/8/2p2P1P/R2b2K1/8 w - - 9 57 | Rc2
Watch how miserable the rook is!

@ 8/8/8/5p2/5k2/2p2P1P/2Rb2K1/8 w - - 23 64 | Ra2?
White must never allow Black to play ...\K e3. 64.\K f2.

% Diagram 209 of the 1500-1599 book: Tsarouchas - Ieromonachos
@ 8/8/8/1p4R1/pK6/8/7p/5b1k w - - 2 68 | Rg4?
White is three pawns down, but he can still hold. First, one must reckon with Black's main threat:
...\B g2, then ...\K g1 and ...h1=\Q. Second, one must realize that when Black plays ...\B g2, his
b5-pawn is unprotected. Third, one must realize that if White can remove the b5-pawn and the
h-pawn, the remaining ending of the \term{wrong-bishop} and rook's pawn is a
draw.\index{Theoretical draws!Wrong bishop and rook's pawn against king} Putting these three
together, White finds the only move that saves the game, like a miracle:\index{Study-like positions}
68.\K a3 \B g2 69.\R xb5 \K g1 70.\R h5.

% Diagram 210 of the 1500-1599 book: Gurung - Mahadevan
@ 8/4k3/p4r2/P1K5/1PB5/8/2P5/8 b - - 4 59 | Rd6?
Black must reckon that White intends to play b5. He must bring his king over to defend against the
passed pawns, and there is no time to lose.\index{Pawn races!Do not lose a tempo} 59...\K d7 60.b5
(60.\B d5 \K c7 61.c4 \K b8 62.\B c6 \K a7 63.\B d5) 60...axb5 61.\B xb5+ \K c7 62.\K b4.

@ 8/4k3/p2r4/P1K5/1PB5/8/2P5/8 w - - 5 60 | Bd5?
There is no time to lose: this is also true for White. 60.b5 axb5 61.\B xb5.

@ 8/4k3/p2r4/P1KB4/1P6/8/2P5/8 b - - 6 60 | Rf6
60...\K d7 61.c4 \R h6 62.b5 axb5 63.cxb5 \K c7. White threatens to play \B c6, an
\term{interference}.\index{Bishop!Interference} As in the bishop endings, we must use our king to
prevent it. 64.a6 \K b8, and Black still holds.

@ 8/4k3/p4r2/P1KB4/1P6/8/2P5/8 w - - 7 61 | c4
Now the pawns are unstoppable.

% Diagram 211 of the 1500-1599 book: De Leon - Manteiga
@ 2k5/8/2P4p/5K1P/5BP1/8/3r4/8 b - - 6 59 | Rc2?
Black's rook can \vocab[strategy]{tie-down}{tie down}{Attack a weakness, or pin a piece, so that an
enemy piece must stay where it is to defend. A tied-down piece cannot take part in the play, and
when it leaves, the weakness falls.}\index{Strategy!Tying down} the white king. White's weakness is
the g4-pawn, which only the king protects. After 59...\R f2 the bishop is pinned, and White cannot
get rid of the pin: if the king runs away, ...\R g2 attacks the pawn, and the king must go back.
59...\R f2 60.\K e5 \R g2 61.\K f5 \R f2. When White gives up the g-pawn, the rook simply ties the
king to the h-pawn instead: 62.c7 \R f1 63.g5 hxg5 64.\K xg5 \R g1+ 65.\K f5 \R h1 66.\K g4 \R g1+
67.\K f5 \R h1, and White makes no progress.

@ 2k5/8/2P4p/5K1P/5BP1/8/2r5/8 w - - 7 60 | c7
60.\B xh6 \R xc6 61.\B g7 \R c5+ 62.\K g6, and Black cannot stop the pawns.

@ 8/2Pk4/6Kp/7P/5BP1/8/2r5/8 b - - 2 61 | Rc6+
61...\R g2, and the rook ties the king down again!\index{Strategy!Tying down}

@ 4k3/2P3K1/2r4p/7P/5BP1/8/8/8 w - - 5 63 | Bxh6?
Exchanging pawns only helps the defender.\index{Defensive resources!Exchanging or removing the pawns}
63.\B g3 \K d7 64.g5! is a \term{breakthrough}:\index{Pawn!Breakthrough} after 64...hxg5 65.h6
the h-pawn runs, and the black king is too far away to stop it.

% Diagram 212 of the 1500-1599 book: Coimbra - Oliveira
@ 8/8/4p3/3p4/3P1kp1/8/4K2b/1R6 w - - 0 55 | Kd3?
Black threatens to play ...\K e4, hitting the weak d4-pawn, and also ...g3. White tries to parry
the threats one by one (55.\K f2 \K e4 56.\R e1+ \K f5 57.\K g2, and now both threats are parried).
It also helps to realize that Black has a weakness of his own on e6. 55.\R b8 \K e4 56.\R e8
\K xd4 57.\R xe6 g3 58.\K f3 \K c4 59.\K g2 \K b3. There are certainly more ways to hold, but we
would like to reach the Ponziani position: 60.\R d6 \K c4 61.\R xd5 \K xd5, and now we have the
\vocab[bishop endings]{ponziani}{Ponziani position}{The defending king stands in front of a rook's or
knight's pawn, and the attacker's bishop is shut in the corner behind its own pawn, for example
bishop on h2 and pawn on g3 against a king on g2. The bishop can never get out, and the game is a
draw.}.\index{Theoretical draws!Ponziani position}
\keysquares{8/8/8/3k4/8/6p1/6Kb/8 w - - 0 62}{}{The Ponziani position: the bishop cannot escape from its corner.}
